\section{Case: chain of thought and proof kernels}
\label{sec:reasoning}

\paragraph{Background.} A large language model computes its output in a \emph{forward
pass}: a fixed sequence of arithmetic operations on its weights and on internal values
called activations. Many models are trained or prompted to write intermediate steps
before the final answer, a \emph{chain of thought}; each written token is fed back as
input to the next forward pass. Where the task is a formal proof, the written steps can
be a proof script that a proof assistant checks, as in recent systems for competition
mathematics \citep{hubert2025}. An explanation is called \emph{faithful} when it
accurately describes the process that produced the output, as opposed to \emph{plausible}
when it is convincing to a reader \citep{jacovi2020}.

\subsection{One word, five layers}

``The model reasons'' is used for five things, which lie in different structures and are
linked by edges of different types (Table~\ref{tab:reasoning}, Figure~\ref{fig:episode}).
By predicate locality (A2) the word may be reused across them only where its truth
conditions stay the same, and they do not.

\begin{table}[t]
\centering\small
\begin{tabularx}{\textwidth}{@{}P{1.8cm}P{5.2cm}P{2.4cm}L@{}}
\toprule
Layer & What ``reasoning'' means there & Structure & Edge and its type \\
\midrule
Activations &
  The forward pass as arithmetic. &
  S1 substrate &
  Constitutive: the computation is realised by the activations; no time elapses. \\
Mechanism &
  Features and circuits across the network's layers, for example an internal two-step
  lookup found in attribution graphs \citep{lindsey2025}. &
  S2 events &
  Causal: the mechanism produces the text; testable by intervention. \\
Trace &
  The chain of thought, or the proof script, that the model writes. &
  S3 representation &
  Representational, back to the mechanism: the text describes how the answer was
  reached, and can misdescribe it. Also causal, forward: the written tokens are input
  to later computation. \\
Validity &
  Each written step follows from the previous ones. &
  S3 under a rule scheme &
  Classificatory: the trace satisfies declared rules. A proof kernel checks this. \\
Access &
  Probes, attribution graphs, reading the trace. &
  S4 epistemic &
  Evidential: an observation supports an inference about the model. \\
\bottomrule
\end{tabularx}
\caption{Five things called ``reasoning'' in a language model.}
\label{tab:reasoning}
\end{table}

\begin{figure}[t]
\centering
\begin{tikzpicture}[
    font=\small,
    box/.style={draw, rounded corners=2pt, align=center, minimum width=6.0cm,
                minimum height=1.0cm, inner sep=4pt},
    side/.style={draw, rounded corners=2pt, align=center, minimum width=2.6cm,
                minimum height=1.0cm, inner sep=4pt},
    lab/.style={font=\footnotesize, align=center},
    >={Stealth[length=2.2mm]}
  ]
  \node[box] (act)   at (0,0)    {S1 \textperiodcentered{} Activations\\\footnotesize weights, arithmetic};
  \node[box] (mech)  at (0,3.1)  {S2 \textperiodcentered{} Mechanism\\\footnotesize features, circuits};
  \node[box] (trace) at (0,6.2)  {S3 \textperiodcentered{} Trace\\\footnotesize chain of thought, proof script};
  \node[side] (rules) at (6.5,6.2) {Rules\\\footnotesize declared scheme};
  \node[side] (obs)   at (-6.5,3.1) {S4 \textperiodcentered{} Observer\\\footnotesize interventions};

  \draw[->] (act) -- node[lab, right] {constitutive\\(no $\Delta t$)} (mech);
  \draw[->] ([xshift=-22mm]mech.north) -- node[lab, left] {causal:\\produces} ([xshift=-22mm]trace.south);
  \draw[->, dashed] (trace.south) -- node[lab, right] {representational:\\faithful?} (mech.north);
  \draw[->] ([xshift=27mm]trace.south) -- node[lab, right] {causal:\\read back} ([xshift=27mm]mech.north);
  \draw[->, very thick] (trace) -- node[lab, above] {classificatory:\\valid} (rules);
  \draw[->] ([yshift=2.5mm]obs.east) -- node[lab, above] {intervention\\(causal)} ([yshift=2.5mm]mech.west);
  \draw[->, dotted, thick] ([yshift=-2.5mm]mech.west) -- node[lab, below] {evidential} ([yshift=-2.5mm]obs.east);
\end{tikzpicture}
\caption{One reasoning episode, with typed edges. A check of validity (thick) runs from
the trace to the rules. The edge that says whether the trace describes the mechanism
(dashed) runs from the trace to the mechanism. The same text is also read back as input
to later computation, a causal edge of its own. The observer intervenes, on the
mechanism as drawn or on the text that is read back; what comes back is evidence
(dotted).}
\label{fig:episode}
\end{figure}

None of the five is the real reasoning with the others as metaphors. Each is a legitimate
object with its own truth conditions. The difficulty arises only when a finding about one
is reported as a finding about another: when a valid proof is taken as evidence of how
the model computed, or a fluent chain of thought as a view of the mechanism.

\subsection{The trace is a representation}

A chain of thought is text the model writes. Its relation to the computation that
produced it is representational: the text describes a route to the answer. Like any
representation it can fail by misrepresenting. \citet{turpin2023} biased models toward
particular answers, for instance by making the correct option always ``(A)'' in the
examples, and found that the models followed the bias while their written explanations
did not mention it. \citet{chen2025} found the same pattern in models trained to reason
before answering: given a hint that changed their answer, they often did not say so in
their chain of thought.

Such findings are often reported with outcome labels: the model ``lied'' or
``rationalised''. Each label adds a claim about intent or mechanism that the observation
does not support. The observation is narrower: a representational edge failed. The text
describes one route; the computation took another. Naming the edge type names the
failure condition, correspondence, and with it what a test has to compare: the text
with the process that produced it.

The same text has a second role. Each written token is input to later forward passes,
so the trace is also a cause of later computation. The two roles are edges of different
types and can come apart: a text can influence the answer without describing how the
answer was reached, and it can describe without influencing. \citet{lanham2023} test the
causal role. They perturb the written reasoning, by truncating it or inserting
mistakes, and measure how far the final answer changes, and they find that the
dependence varies with model and task.

\subsection{What a proof kernel certifies}
\label{sec:kernel}

Now replace the chain of thought by a proof in Lean. The elaborator turns what the model
wrote into a proof term, and the kernel checks that term against the stated proposition,
in a declared environment of definitions and axioms. If the kernel accepts, the
proposition is proved relative to that environment, a guarantee no informal chain of
thought has. The edge that was checked is the classificatory edge from the trace to the
rules. The representational edge from the trace to the mechanism was not in view.

This can be stated exactly in a small model.

\begin{definition}[Reasoning system]
\label{def:rs}
A \emph{reasoning system} consists of a set $M$ of mechanisms, a set $T$ of traces, a
function $\produce : M \to T$, a predicate $\valid \subseteq T$ and a relation
$\faithful \subseteq T \times M$.
\end{definition}

\noindent
Here $\produce$ stands for the causal edge from mechanism to trace, $\valid$ for the
classificatory edge from trace to rules, and $\faithful(t,m)$ for the representational
edge: the trace $t$ describes the mechanism $m$. Write $F(m)$ for
$\faithful(\produce(m), m)$: the trace that $m$ writes describes $m$. A \emph{judgement
on traces} is a predicate $D \subseteq T$. It \emph{agrees with faithfulness} when, for
all $m \in M$, $D(\produce(m))$ holds if and only if $F(m)$.

\begin{proposition}[Validity is blind to the mechanism]
\label{prop:blind}
If $\produce(m_1) = \produce(m_2)$, then $\valid(\produce(m_1))$ if and only if
$\valid(\produce(m_2))$.
\end{proposition}

\begin{proposition}[When a judgement on traces can agree with faithfulness]
\label{prop:char}
A judgement on traces that agrees with faithfulness exists if and only if faithfulness
is constant among the mechanisms that write the same trace: for all $m_1, m_2 \in M$,
if $\produce(m_1) = \produce(m_2)$ then $F(m_1)$ if and only if $F(m_2)$.
\end{proposition}

\begin{corollary}[The obstruction]
\label{cor:split}
If two mechanisms write the same trace and the trace describes one of them and not the
other, that is, $\produce(m_1) = \produce(m_2)$, $F(m_1)$ and not $F(m_2)$, then no
judgement on traces agrees with faithfulness.
\end{corollary}

\noindent
These results need no assumption about the relation $\faithful$. They are checked in
Lean (Section~\ref{sec:mech}); the proofs take a few lines each. A judgement on traces
can be anything that looks only at the trace: a proof kernel, a reward model, a human
grader, a comparison with a reference answer. Suppose two routes, one following the
stated steps and one using an unmentioned hint, write the same valid text, and the text
describes the first and not the second. Then no judgement on the text sorts the
faithful episode from the unfaithful one. Appendix~\ref{app:uniqueness} gives one
condition under which a shared trace must have different verdicts, and the two-route
example from the package.

\begin{remark}[The results are elementary]
Proposition~\ref{prop:blind} is the congruence of equality.
Proposition~\ref{prop:char} says that a function of the trace can track a property of
mechanisms exactly when that property does not vary among the mechanisms behind one
trace. The paper claims no depth for these statements. Their use is that they fix which
edge a check tests, in a setting where one word lets a guarantee about one edge pass for
a guarantee about another.
\end{remark}

\begin{remark}[Scope]
\label{rem:scope}
(a) \emph{Existence, not computation.} A judgement on traces is an arbitrary predicate.
Proposition~\ref{prop:char} says when one exists that agrees with faithfulness, not
that it can be computed or found. In particular, a shared trace is by itself no
obstruction: if two mechanisms write the same trace and it describes neither, the
judgement that rejects the trace is right about both.
(b) \emph{Decision, not evidence.} The results do not say that a trace carries no
evidence about the mechanism. A feature of the text may correlate with faithfulness
across many episodes. What Corollary~\ref{cor:split} excludes is a judgement on the
text that is right on both episodes of a shared trace with different verdicts.
(c) \emph{Production as a function.} $\produce$ is a function for simplicity. The
argument uses only that one trace occurs with two mechanisms.
\end{remark}

Whether real models contain two computations that write an identical trace with
different verdicts is an empirical question that this paper does not answer. The
studies cited above show episodes in which the written reasoning omits a factor that
changed the answer. They motivate the distinction between validity and faithfulness.
They do not exhibit such a pair, and were not designed to.

\subsection{Two clocks: layers and tokens}

Within one forward pass, a computation can only be as serial as the network is deep: each
layer is one step, and the depth is fixed. A chain of thought adds a second sequence,
token positions. Every written token triggers a new forward pass that can read
everything before it, so the serial depth available grows with the length of the trace.
\citet{merrill2024} make this precise for a class of transformers. With logarithmically
many intermediate steps the gain is slight; with linearly many, under an architectural
assumption, all regular languages can be recognised; with polynomially many, the class
is exactly the problems solvable in polynomial time. Under standard
complexity-theoretic assumptions these are strict extensions of what the same
transformers compute without intermediate steps.

So the question of how many steps a model can reason has no answer until the sequence is
declared: layers or tokens. The two lie in different places in
Table~\ref{tab:reasoning}. Depth in layers is inside the mechanism (S2). Depth in tokens
runs through the trace (S3), because the intermediate results are written down and read
back. Approaches that feed continuous hidden states back in place of tokens, such as
that of \citet{hao2024}, keep the second sequence and remove the written layer. The
computation is no less real for that. It is less legible: the representation that could
at least be read, and sometimes checked, is gone.

\subsection{Grading the answer, grading the steps}

Training follows the same division. Outcome-based rewards grade only the final answer.
Process supervision grades each step, and on competition mathematics it trained more
reliable reward models than outcome supervision did \citep{lightman2023}. A graded step
is still a written step. Process supervision moves the judgement from the end of the
trace to its inside; it does not move it off the trace's layer, and
Proposition~\ref{prop:char} applies to it unchanged.

What goes beyond a judgement on a fixed text is an intervention, and there are two
places to intervene. One is the text as input: the written reasoning is perturbed and
the change in the answer observed, as in the study of \citet{lanham2023}. This needs no
access to the model's internals. The other is the mechanism: ablation and activation
patching, and attribution graphs that trace which internal features drive which,
validated by such interventions \citep{lindsey2025}. In both cases the observer
obtains, by an evidential edge, an observation of how the process behaves under a
change, and not only an observation of the text. These methods are costly and partial.
A claim that a trace is faithful is a claim about a representational edge. Where
faithfulness differs between computations that write the same trace, an exact judgement
needs information beyond the trace (Corollary~\ref{cor:split}), and interventions are
one way to obtain it.

\paragraph{What the field does, and what the rule adds.} Interpretability research
already separates faithfulness from plausibility and already tests explanations by
intervention; causal abstraction gives that practice a formal theory
(Section~\ref{sec:related}). The rule adds a plain statement of what follows for checks
on text. Validity is a classificatory edge from text to rules. Faithfulness is a
representational edge from text to mechanism. Whatever a check on the trace establishes,
it establishes about the trace, and a model's own account of how it reasoned is one more
trace. This does not make verified proofs less valuable: a valid proof is valid, whatever
produced it. It limits what else may be concluded from the check.
